\begin{figure}[tbp]
\centering
\begin{tikzpicture}[x=1cm,y=1cm,font=\small,
  dot/.style={circle,fill,inner sep=1.6pt},
  raw/.style={very thick,blue!65!black,-{Stealth[length=2mm]}},
  escape/.style={thick,orange!70!black,-{Stealth[length=2mm]}}]
\node[dot,label={above left:$u=g_0$}] (u) at (0,0) {};
\node[dot,label=above:$g_J$] (a) at (1.5,1.35) {};
\node[dot,label=above:$h_J$] (b) at (6.5,1.35) {};
\node[dot,label={above right:$v=h_0$}] (v) at (8,0) {};
\node[dot,label=left:$g_{-J}$] (c) at (-1.1,-1.1) {};
\node[dot,label=right:$h_{-J}$] (d) at (9.1,-1.1) {};
\node[dot,label=below:$z$] (z) at (0,-2.65) {};
\node[dot,label=below:$z'$] (zz) at (8,-2.65) {};
\draw[raw] (u)--node[above left] {$J$} (a);
\draw[raw] (a) to[out=-10,in=170] (2.6,1.1)
  to[out=-10,in=180] (4,1.3) to[out=0,in=190] (b);
\node[align=center] at (4,2.65) {middle walk of length $L$\\
conditioned to end in the type of $g_J$};
\draw[raw] (b)--node[above right] {$J$} (v);
\draw[escape] (u)--node[left] {$J$} (c);
\draw[escape] (v)--node[right] {$J$} (d);
\draw[escape] (c) to[out=-80,in=180] (-.3,-1.9) to[out=0,in=180] (.5,-2.1)
  to[out=0,in=0] (z);
\draw[escape] (d) to[out=-100,in=0] (8.3,-1.9) to[out=180,in=0] (7.5,-2.1)
  to[out=180,in=180] (zz);
\node[align=center,text width=5cm] at (4,-1.2)
{Walk tails\\of length $m=j-J$};
\end{tikzpicture}
\caption{The anchored path experiment. The blue raw path has a
prescribed prefix, a reversible middle walk, and a prescribed suffix.
The two orange escapes use the opposite halves of the same geodesic
segments and walk tails that, conditional on $g,\eta,h$, are independent
ordinary $P$-walks before avoidance is imposed. Reversing the middle
walk and exchanging the two segments and escape tails gives the reversed
raw-path law after transport on the center orbit. Avoidance is tested
against the raw path; the escapes may intersect each other. Loop-erasing
either escape keeps it outside the corridor interior.}
\label{fig:escape-templates}
\end{figure}
